\begin{abstract}
Modern GPU-centric systems run heterogeneous workloads whose performance and efficiency largely depend on policies for memory placement, kernel scheduling, and observability. Today, these policies are either implemented inside framework-specific runtimes or hard-coded into proprietary vendor stacks, without programmable OS-level interface. Consequently, adapting to changing workloads, SLAs, or hardware still requires invasive code changes and operational disruptions. We argue that the OS GPU driver and device layer, uniquely positioned with cross-tenant visibility and fine-grained control over privileged mechanisms, must serve as a programmable policy narrow waist. However, extending CPU-side programmability frameworks like eBPF to GPUs is challenging: GPU drivers tightly couple resource management with low-level hardware interactions, and host-only or device-only approaches fail to deliver coherent, cross-layer policies.

We present \sys, an eBPF-based policy substrate that spans CPU hosts and GPU devices. \sys refactors GPU drivers to expose safe, verified eBPF hooks at key control points, and offloads these policies onto GPU devices via a lightweight sandboxed JIT, coordinated by cross-layer states and single control plane. This unified runtime allows workload-specific memory offloading, kernel preemption, and observability policies to be expressed and enforced consistently across kernel and device contexts, providing strong safety guarantees and low overhead. We implement \sys by extending Linux GPU drivers and building a GPU-side runtime. Evaluation on realistic workloads, including inference, training, and vector search, shows that \sys-based policies improve throughput by up to $5\times$, reduce tail latency by up to $2\times$, and incur only a few percent observability overhead. These results demonstrate the practicality of a driver-anchored, programmable CPU–GPU policy substrate.

\xiangyu{I think it is better to reduce the abstract content by 50-60\% by removing some example details.}

\xiangyu{Here is my version
The rapid development of various LLM applications makes GPU-centric systems run heterogeneous workloads. However, recent vendors offer hard-coded proprietary resource management and allocation policies without enabling users to easily implement application-specific ones. This has a negative effect on the final execution performance and efficiency.

We argue that these policies should be programmable to offer users fine-grained control and build gBPF, an eBPF-based policy substrate that spans CPU hosts and GPU devices. gBPF extends eBPF to GPU by extending Linux GPU drivers and building a GPU-side runtime. Results show that gBPF offers GPU-side programmability with safety guarantees and low overhead. gBPF-based policies can improve throughputs by up to 5x and reduce tail latency by up to 2x on realistic workloads.}
\end{abstract}

\section{Introduction}

GPUs now dominate the compute and capital costs of modern computing systems, powering workloads that range from latency-sensitive LLM inference pipelines to large-scale training, graph analytics, and vector search. Across these workloads, end-to-end performance and efficiency depend not just on raw FLOPs, but on \emph{policies}: how the system places and migrates data across device and host memory, admits and schedules kernels in multi-tenant environments, and collects observability signals for online adaptation. These policies must evolve continuously as models, SLAs, and hardware change. 

Yet today, these policies, offered by vendors, remain largely fixed within proprietary stacks and framework-specific runtimes, making GPU infrastructure rigid and costly to operate. 
% In current GPU stacks, the host driver and device firmware implement fixed, opaque mechanisms for memory management and offloading, context scheduling, and performance event collection. 
Frameworks such as PyTorch or TensorFlow, as well as emerging inference runtimes like vLLM, build bespoke management logic atop these mechanisms, often defensively preallocating GPU memory or embedding scheduling heuristics tailored to a specific workload mix, and in some cases exposing framework-local policy APIs. 
Since these underlying driver and firmware mechanisms expose no safe extension points, operators cannot implement cross-framework, workload-specific policies for memory placement or preemption precisely.
% where decisions occur, such as during device memory migration or when low-priority tasks risk delaying latency-critical queries. 
Consequently, GPU deployments frequently suffer from unpredictable tail latency, underutilization, and operational disruptions whenever policies require updates or fine-tuning. 
User-space sharing and serving frameworks can therefore be programmable \xiangyu{in a restricted way to the behavior that the fixed driver and device mechanisms choose to expose.}
% , but their control is necessarily limited to 

% \xiangyu{This paragraph shows the process of our design, I think we should move it to later part.}
% By contrast, modern CPU systems have transitioned toward programmable policies. The extended Berkeley Packet Filter (eBPF) ecosystem supports dynamically loaded, verified programs for scheduling (\texttt{sched\_ext}), page-cache eviction (\texttt{cache\_ext}), TCP congestion control (\texttt{struct\_ops}), and rich observability~\cite{gregg16,gregg2021computing,Bachl21,Zhong22,jia2023programmable}. These designs expose stable kernel mechanisms through structured attach points, enabling dynamic policy updates without kernel reboots or patches~\cite{schedext-docs,cacheext-docs}. Applying a similar approach to GPUs and treating the GPU driver as another kernel subsystem and attaching host-side eBPF programs is appealing. Conversely, GPU-only instrumentation frameworks such as CUPTI, NVBit, eGPU, and Neutrino~\cite{cupti,villa2019nvbit,neutrino,egpu} provide telemetry via binary rewriting at the PTX or SASS level, but primarily target profiling, lack unified safety guarantees, and do not integrate with privileged host mechanisms.

Our observation is that GPU applications need an OS-level narrow waist for policy: the GPU driver and device execution environment are uniquely positioned with global cross-tenant visibility and direct control over privileged mechanisms such as page migration, kernel scheduling, and preemption. We therefore treat GPU resource management as a first-class kernel subsystem and make it programmable by exposing a small set of verifiable policy interfaces.
The right abstraction for GPU-centric programmability treats CPU hosts and GPU devices as a single coherent policy domain, supported by a unified intermediate representation, shared state abstraction, and a common control plane with lifecycle management. 

\xiangyu{The recent development of the extended Berkeley Packet Filter (eBPF) ecosystem in CPU systems gives a possibility to support programmability in the GPU domain. 
Concretely, eBPF’s restricted instruction set, verifier guarantees, and rich existing ecosystem uniquely enable such a unified GPU policy substrate, seamlessly extending unified programmability from CPUs to GPUs.}

However, directly applying host-only eBPF to GPUs fails for two fundamental reasons. 
First, critical decisions such as thread-block scheduling, warp prioritization, or identifying pathological CUDA kernel behaviors occur device-side within GPU kernels and firmware, invisible to host-only instrumentation. 
Second, one-size-fits-all GPU drivers tightly couple resource management with hardware interactions (e.g., MMU operations, page fault handling, context management) without stable, safe hooks; naive insertion of probes risks GPU hangs, kernel panics, or resource leaks. 
A practical GPU policy substrate must therefore separate mechanism from policy inside performance-critical GPU drivers without destabilizing hardware interactions, and safely execute policy logic within massively parallel, performance-sensitive GPU kernels.

\xiangyu{There are multiple challenges to designing a general substrate for GPU policy programmability.}
% Designing a general substrate for GPU policy programmability therefore introduces several challenges. 
\emph{C1: Safe and efficient driver extensibility.} We must retrofit mechanism-policy separation into large, performance-critical GPU drivers without destabilizing their low-level interactions with hardware. \emph{C2: Safe and efficient device-side execution.} Policy logic must execute directly within GPU kernels on critical memory access and synchronization paths, while ensuring bounded execution and memory safety in a massively parallel environment. \emph{C3: Cross-layer state and abstractions.} Policies often require correlating host-side and device-side signals (e.g., page access patterns with warp-level stalls), necessitating a unified state abstraction such as time and virtual address, and efficient synchronization across the CPU-GPU boundary. \emph{C4: A single control plane.} The policy runtime must support unified verification, deployment, and debugging across both host and device execution contexts to avoid introducing additional fragmentation.

We present \sys, a unified, safe, and efficient policy runtime for GPU-centric systems designed explicitly to address these challenges. \sys exposes a small, stable set of GPU mechanisms for memory placement and migration, kernel admission and preemption, and kernel-level instrumentation as programmable eBPF attach points within an extended GPU driver. Policies are written as eBPF programs using a restricted instruction subset, statically verified for memory safety and bounded execution. The identical eBPF intermediate representation is then offloaded onto GPU devices: \sys compiles verified bytecode into a GPU-specific device code representation, injects it via lightweight trampolines at kernel launch time, and executes it within a sandboxed GPU-side runtime. Both host-side and device-side policy programs share a common map abstraction backed by unified memory, enabling low-latency, coherent state sharing without explicit data marshalling. A single verifier and control plane manage policy deployment and lifecycle across host and device, forming a unified policy domain.

We implement \sys on Linux by extending \sysdriver using NVIDIA's open GPU kernel modules~\cite{}, constructing a GPU-side device JIT and sandbox runtime. Driver modifications adopt a \texttt{struct\_ops}-style approach, introducing carefully chosen attach points within memory management and scheduling paths, delegating policy decisions to verified eBPF callbacks. On the GPU device, lightweight trampolines save minimal register state, invoke translated eBPF snippets, and restore state, introducing sub-microsecond overhead. Unlike prior GPU observability tools~\cite{cupti,villa2019nvbit,neutrino,egpu}, \sys provides a single unified, framework-agnostic policy abstraction across driver and device contexts under one verifier and control plane.

We evaluate \sys across realistic GPU workloads including LLM inference pipelines, GNN training, vector search, and mixed-priority multi-tenant scenarios. Using \sys, we implement adaptive memory prefetching and eviction policies, fine-grained kernel preemption strategies, and bcc-style GPU kernel observability tools. \sys-based policies improve throughput by up to $5\times$ and reduce tail latency by up to $2\times$ compared to static heuristics, while instrumentation-only deployments incur only $2$–$3\%$ overhead. 
% These results confirm that unified, eBPF-based policy programmability spanning CPU and GPU contexts represents a practical and powerful abstraction for managing GPU-centric systems.
\xiangyu{Our contributions can be summarized as following:}

\begin{itemize}[leftmargin=*,itemsep=0pt]
    \item We identify 
    % and characterize key 
    programmability gaps in GPU stacks for resource management and scheduling, and show how the lack of a programmable, kernel-anchored policy waist at the GPU device layer makes it difficult to adapt to \xiangyu{customized requirements}.
    \item We implement \sys, an eBPF-based policy runtime that treats the GPU driver and device execution layer as a unified OS subsystem. 
    % \sys exposes verified driver hooks for memory and scheduling, executes the same eBPF IR inside GPU kernels via a sandboxed device-side JIT, and shares state through unified-memory-backed BPF maps under a single verifier and control plane.
    \item We demonstrate that \sys enables adaptive, workload-specific policies that significantly improve performance
    % , resource utilization, and tail latency 
    on realistic LLM 
    % inference, training, and vector 
    workloads, with low overhead.
    % and without requiring application or driver restarts.
\end{itemize}
